\documentclass[sigconf,nonacm]{acmart}

\usepackage{revision-marks}
\showrevmarksfalse

\title{Selecting Latency Budgets for Edge Inference}
\author{San Zhang}
\affiliation{%
  \institution{University of Technology}
  \city{Example City}
  \country{China}%
}

\begin{document}
\begin{abstract}
We compare two inference schedules under a fixed latency budget. This example
uses request traces with different arrival patterns to isolate the effect of
queueing on deadline misses.
\end{abstract}
\maketitle

\section{Introduction}
On-device inference services must allocate memory and compute before a request
arrives. A schedule that performs well on average can still exceed the deadline
when the request queue grows. We study a schedule that reserves capacity for
these bursts and compare it with a fixed-batch baseline.

\section{Measurement Protocol}
\rpoint[E3]We measure latency from queue entry to result delivery, including
preprocessing and device transfers. Each run starts with a fixed warm-up period
followed by the same sequence of requests. We replay the sequence for each
schedule so the comparison uses the same arrival times.

% Keep the four columns stable across TeX Live versions.
\newpage
\section{Baseline Schedule}
A fixed-batch scheduler accumulates requests until the batch is full or its
timer expires. We use the same timer in each run and size the batch using a
separate calibration trace. This rule provides a reproducible baseline for
assessing deadline misses.

\rpoint[E4]The calibration trace is disjoint from the evaluation trace. Once a
batch size is selected, it is held fixed for all reported runs.

\clearpage
\section{Evaluation}
We evaluate the schedules on a steady workload and a bursty workload. The
steady workload has uniformly spaced arrivals. The bursty workload preserves
the same total request count while changing the arrival times.

We report the median and 95th percentile end-to-end latency for each workload,
along with the fraction of requests that miss the deadline. The fraction is
calculated over all submitted requests, including those left in the queue at
the deadline.

\newpage
\section{Budget Selection}
The adaptive schedule maintains a running estimate of the time needed to
complete the current batch. A new request joins the batch when it can be served
within the remaining budget.

Each decision uses only the current queue and the completion times of earlier
batches. We initialize the estimate from the calibration trace and update it
after each batch completes. When the remaining budget is too short to add the
next request, the scheduler dispatches the batch immediately. The pending
request then starts a new batch with its original queue entry timestamp.

\rpoint[E1]\rpoint[E2]We set the budget to 40\,ms before evaluating either
workload and use the same value for every device. This budget includes time
spent in the queue and time spent on device.

\section{Conclusion}
A shared budget and a consistent latency definition make it possible to
compare the schedules across arrival patterns and devices.
\end{document}
